% Intensité relative de la lumière diffusée par l'air selon la longueur d'onde (corrigé du défi D19)
\begin{tikzpicture}
\begin{axis}[ybar,bar shift=0pt,bar width=7mm,width=100mm,height=56mm,ymin=0,ymax=10.5,
  xtick={400,450,...,700},xmin=370,xmax=730,
  xticklabel style={font=\otLight\fontsize{7}{8}\selectfont,text=otMuted},
  yticklabel style={font=\otLight\fontsize{7}{8}\selectfont,text=otMuted},
  axis lines=left,axis line style={otMuted,line width=0.5pt},tick style={otMuted},
  xlabel={$\lambda$ (nm)},ylabel={diffusion relative (\qty{700}{nm} : 1)},label style={font=\small,text=otInk},clip=false]
\addplot[fill=violet!70,draw=none] coordinates {(400,9.38)};
\addplot[fill=blue!70,draw=none] coordinates {(450,5.86)};
\addplot[fill=cyan!70!blue,draw=none] coordinates {(500,3.84)};
\addplot[fill=green!65!black,draw=none] coordinates {(550,2.62)};
\addplot[fill=yellow!80!orange,draw=none] coordinates {(600,1.85)};
\addplot[fill=orange,draw=none] coordinates {(650,1.35)};
\addplot[fill=red!80,draw=none] coordinates {(700,1.0)};
\end{axis}
\end{tikzpicture}
\par\smallskip\legende{La diffusion varie comme $\dfrac{1}{\lambda^4}$ : le bleu est diffusé près de six fois plus que le rouge.}
